% ===========================================================================
%  The entanglement programme: quantum information from the shadow tower
% ===========================================================================
%
%  This section develops the systematic connection between the proved
%  algebraic core (shadow tower, bar-cobar, complementarity, open/closed)
%  and the entanglement programme in quantum gravity (RT, QES, islands,
%  Page curve, error correction, modular flow).  The unique asset: the
%  shadow tower provides a constructive, convergent, finite-order
%  approximation scheme — contrasting with the asymptotic genus
%  expansion of string theory and the uncontrolled spin-foam sums
%  of loop quantum gravity.
%
%  Status hierarchy:
%    ProvedHere    — follows from the existing proved core
%    Conjectured   — well-motivated, precise, falsifiable
%    Heuristic     — structural identification, not yet a theorem

\section{The entanglement programme: quantum information from the shadow tower}
\label{sec:thqg-entanglement-programme}
\index{entanglement programme|textbf}
\index{quantum information!from shadow tower|textbf}

% ======================================================================
%
%  1. THE ENTANGLEMENT DICTIONARY
%
% ======================================================================

\subsection{The entanglement dictionary}
\label{subsec:thqg-entanglement-dictionary}
\index{entanglement dictionary|textbf}

The holographic modular Koszul datum
$\cH(\cA) = (\cA, \cA^!, \mathcal{C}, r(z), \Theta_\cA,
\nabla^{\mathrm{hol}})$
(Definition~\ref{def:thqg-holographic-datum}) encodes six algebraic
structures.  Each has a precise counterpart in the holographic
entanglement programme:

\begin{center}
\renewcommand{\arraystretch}{1.4}
\small
\begin{tabular}{l@{\;\;}l@{\;\;}l@{\;\;}c}
\textbf{Algebraic object} & \textbf{Holographic role}
  & \textbf{Entanglement concept} & \textbf{Status} \\
\midrule
$Q_g(\cA) \oplus Q_g(\cA^!)$
  & Lagrangian splitting of phase space
  & bipartite decomposition
  & \ref{thm:thqg-g3-polarization} \\
$\sigma$ (Verdier involution)
  & gravitational CPT
  & Tomita--Takesaki $J$
  & \S\ref{subsec:thqg-entanglement-modular-flow} \\
$\langle \cdot, \cdot \rangle_g$ (Verdier pairing)
  & bulk--boundary coupling
  & entanglement bilinear form
  & \ref{thm:thqg-III-genus-1-holographic} \\
$\Theta_\cA^{(g)}$ (genus-$g$ MC element)
  & genus-$g$ gravitational state
  & entanglement carrier
  & \ref{thm:mc2-bar-intrinsic} \\
$\Theta_\cA^{\le r}$ (shadow truncation)
  & $r$-th order QES approximant
  & finite-order entanglement
  & \S\ref{subsec:thqg-entanglement-qes} \\
$\kappa(\cA) \cdot \lambda_g^{\mathrm{FP}}$
  & Bekenstein--Hawking entropy
  & classical (area) term
  & Theorem~D \\
$Q_L$ (shadow metric)
  & modular Hamiltonian
  & entanglement spectrum generator
  & \S\ref{subsec:thqg-entanglement-modular-flow} \\
$\nabla^{\mathrm{sh}}$ (shadow connection)
  & modular flow
  & entanglement flow
  & Theorem~\ref{thm:shadow-connection} \\
$B(\cA) \dashv \Omega$ (bar-cobar)
  & encoding/decoding
  & error correction
  & \S\ref{subsec:thqg-entanglement-error-correction} \\
Koszulness of $\cA$
  & code perfection
  & exact recovery
  & Theorem~B \\
$r_{\max}(\cA)$ (shadow depth)
  & code complexity
  & correction capacity
  & \S\ref{subsec:thqg-entanglement-error-correction} \\
$\Theta^{\mathrm{oc}}_\cA$ (open/closed MC)
  & subregion/bulk coupling
  & entanglement wedge
  & \ref{thm:thqg-oc-mc-equation} \\
$HH_*(\cA)$ (Hochschild homology)
  & annulus trace
  & partial trace
  & \ref{thm:thqg-annulus-trace} \\
$\kappa + \kappa' = \text{const}$
  & complementarity constraint
  & Page bound
  & Theorem~D
\end{tabular}
\end{center}

\begin{remark}[Three distinct entanglement structures]
\label{rem:thqg-three-entanglements}
\index{entanglement!three structures}
Three entanglement structures arise from the algebraic
framework, corresponding to three physical settings:
\begin{enumerate}[label=\textup{(\roman*)}]
\item \emph{Bulk--boundary entanglement}
  (\S\ref{subsec:thqg-entanglement-higher-genus}):
  the correlation between $Q_g(\cA)$ and $Q_g(\cA^!)$
  in the Lagrangian decomposition.  Measured by the
  \emph{entanglement rank} $E_g(\cA)$
  (Definition~\ref{def:thqg-entanglement-rank}).
  This is the bipartite structure at genus $g$ controlling
  the gravitational free energy decomposition.

\item \emph{Spatial entanglement entropy}
  (\S\ref{subsec:thqg-entanglement-qes}): the entropy
  associated to a boundary subregion, controlled by the
  open/closed MC element $\Theta^{\mathrm{oc}}_\cA$
  (Theorem~\ref{thm:thqg-oc-mc-equation}).  This is the
  algebraic incarnation of the Ryu--Takayanagi/QES formula.

\item \emph{Gravitational entropy}
  (\S\ref{subsec:thqg-entanglement-page}): the Bekenstein--Hawking
  entropy and its quantum corrections, controlled by the shadow
  tower generating function.  Already developed in
  \S\ref{subsec:thqg-VIII-btz-entropy}.
\end{enumerate}
The shadow tower $\Theta_\cA^{\le r}$ controls all three:
the scalar level ($r = 2$) gives the classical term in each,
and higher arities give the quantum corrections.  The
convergence theorem (Theorem~\ref{thm:recursive-existence})
guarantees that these corrections are summable.
\end{remark}

% ======================================================================
%
%  2. GRAVITATIONAL ENTANGLEMENT AT HIGHER GENUS
%
% ======================================================================

\subsection{Gravitational entanglement at higher genus}
\label{subsec:thqg-entanglement-higher-genus}
\index{entanglement!higher genus|textbf}

\begin{definition}[Entanglement rank]
\label{def:thqg-entanglement-rank}
\index{entanglement rank|textbf}
Let $(\cA, X, \langle \cdot, \cdot \rangle, k)$ be a
gravitational input.  The \emph{entanglement rank at genus~$g$}
is
\begin{equation}\label{eq:thqg-entanglement-rank}
E_g(\cA)
\;:=\;
\rank\bigl(\Theta_\cA^{(g)}
\bigr|_{Q_g(\cA) \times Q_g(\cA^!)}\bigr),
\end{equation}
where the rank is computed as follows.  The MC element
$\Theta_\cA^{(g)} \in \cH_g(\cA) = Q_g(\cA) \oplus Q_g(\cA^!)$
has components $\Theta_\cA^{(g)} = (\alpha_g, \beta_g)$
in the Lagrangian coordinates.  The Verdier pairing
$\langle \cdot, \cdot \rangle_g$ induces a bilinear form
$\cB_g^{\Theta} \colon Q_g(\cA) \times Q_g(\cA^!) \to \mathbb{C}$
defined by
\begin{equation}\label{eq:thqg-entanglement-bilinear}
\cB_g^{\Theta}(v, w)
\;:=\;
\langle v, \pi_{Q_g(\cA^!)}(\Theta_\cA^{(g)}) \rangle_g
\;+\;
\langle \pi_{Q_g(\cA)}(\Theta_\cA^{(g)}), w \rangle_g.
\end{equation}
The entanglement rank $E_g(\cA)$ is the rank of
$\cB_g^{\Theta}$ as a bilinear form.
\end{definition}

\begin{remark}[Interpretation]
\label{rem:thqg-entanglement-rank-interpretation}
The entanglement rank measures the number of independent channels
through which $Q_g(\cA)$ and $Q_g(\cA^!)$ communicate via the
MC element.  When $E_g = 0$, the genus-$g$ gravitational state
lives entirely in one Lagrangian summand: no bulk--boundary coupling
at genus~$g$.  When $E_g = \min(\dim Q_g(\cA), \dim Q_g(\cA^!))$,
the coupling is maximal.
\end{remark}

\begin{theorem}[Entanglement structure theorem;
\ClaimStatusProvedHere]
\label{thm:thqg-entanglement-structure}
\index{entanglement structure theorem|textbf}
For a gravitational input $(\cA, X, \langle \cdot, \cdot \rangle, k)$:
\begin{enumerate}[label=\textup{(\roman*)}]
\item \textup{(Von Neumann bound.)}
  $E_g(\cA) \le \min\bigl(\dim Q_g(\cA),\, \dim Q_g(\cA^!)\bigr)$.

\item \textup{(Genus-$1$ triviality.)}
  $E_1(\cA) = 1$ and $\dim Q_1(\cA) = \dim Q_1(\cA^!) = 1$.

\item \textup{(Self-dual saturation.)}
  For $\cA \cong \cA^!$: $\dim Q_g(\cA) = \dim Q_g(\cA^!)
  = \tfrac{1}{2} \dim \cH_g(\cA)$ for all $g \ge 1$.

\item \textup{(Shadow tower control.)}
  The entanglement rank decomposes by shadow arity:
  \begin{equation}\label{eq:thqg-entanglement-arity-decomposition}
  E_g(\cA) \;=\; E_g^{\le 2}(\cA) + \delta E_g^{(3)}(\cA)
  + \delta E_g^{(4)}(\cA) + \cdots,
  \end{equation}
  where $E_g^{\le 2}$ is the scalar contribution \textup{(}at most $1$\textup{)}
  and $\delta E_g^{(r)}$ is the rank increment from the arity-$r$ shadow jet.
  The series stabilizes at $r = r_{\max}(\cA)$ for finite shadow depth,
  and converges for $r_{\max} = \infty$ by
  Theorem~\textup{\ref{thm:recursive-existence}}.

\item \textup{(Complementarity constraint.)}
  $E_g(\cA) = E_g(\cA^!)$.
\end{enumerate}
\end{theorem}

\begin{proof}
\emph{Part (i).}  The rank of a bilinear form
$\cB_g^{\Theta} \colon V \times W \to \mathbb{C}$ is at most
$\min(\dim V, \dim W)$.

\emph{Part (ii).}
At genus $1$, $\dim \cH_1(\cA) = 2$
(Theorem~\ref{thm:thqg-III-genus-1-holographic}(i)),
with $\dim Q_1(\cA) = \dim Q_1(\cA^!) = 1$.  The MC element
$\Theta_\cA^{(1)}$ is nonzero (it includes the scalar
$\kappa(\cA) \ne 0$), so $\cB_1^{\Theta}$ has rank exactly~$1$.

\emph{Part (iii).}
Corollary~\ref{cor:thqg-III-dimension-parity}.

\emph{Part (iv).}
The MC element decomposes by the shadow arity filtration
$F^{\le r}$ on $\gAmod$.  The scalar truncation
$\Theta_\cA^{\le 2,(g)}$ has rank at most~$1$
(it is proportional to $\kappa(\cA) \cdot \lambda_g^{\mathrm{FP}}$,
a scalar).  Each higher arity can increase the rank by at most
$\dim(\text{arity-$r$ component})$.  Convergence of the series
follows from Theorem~\ref{thm:recursive-existence}: the arity-$r$
components decay as $O(\rho(\cA)^r)$ where
$\rho(\cA)$ is the shadow growth rate.

\emph{Part (v).}
The complementarity exchange principle
(Corollary~\ref{cor:thqg-III-complementarity-exchange}):
$\cB_g^{\Theta}$ and $\cB_g^{\sigma(\Theta)}$ have the
same rank since $\sigma$ exchanges $Q_g(\cA)$ and
$Q_g(\cA^!)$.
\end{proof}

\begin{corollary}[Entanglement vanishing criterion;
\ClaimStatusProvedHere]
\label{cor:thqg-entanglement-vanishing}
\index{entanglement!vanishing criterion}
If $\cA$ has shadow depth $r_{\max}(\cA) = 2$
\textup{(}class G: Heisenberg, lattice VOAs\textup{)},
then $E_g^{\le 2}(\cA) \le 1$ for all $g \ge 1$:
the bulk--boundary entanglement is at most one-dimensional at every genus.
For classes L, C, M, the higher shadow jets can
increase $E_g$ above $1$ at genus $g \ge 2$.
\end{corollary}

\begin{proof}
For class~G, $\Theta_\cA = \Theta_\cA^{\le 2}$
(the shadow tower terminates at arity~2), so
$\Theta_\cA^{(g)} = \kappa(\cA) \cdot \lambda_g^{\mathrm{FP}}$,
which is a scalar multiple of a fixed class.  The bilinear form
$\cB_g^{\Theta}$ has rank at most~$1$.
\end{proof}

\begin{proposition}[Entanglement entropy from the shadow tower;
\ClaimStatusHeuristic]
\label{prop:thqg-entanglement-entropy-higher}
\index{entanglement entropy!higher genus}
For a \emph{unitary} chiral algebra $\cA$ with positive-definite
Shapovalov form, the Lagrangian decomposition
$\cH_g(\cA) = Q_g(\cA) \oplus Q_g(\cA^!)$ upgrades to
a Hilbert space tensor product via the unitarity inner product.
In this setting, the MC element $\Theta_\cA^{(g)}$ determines
a reduced density matrix
\begin{equation}\label{eq:thqg-density-matrix}
\rho_g(\cA)
\;:=\;
\Tr_{Q_g(\cA^!)}\bigl(
|\Theta_\cA^{(g)}\rangle \langle \Theta_\cA^{(g)}|
\bigr)
\end{equation}
and the \emph{genus-$g$ entanglement entropy}
\begin{equation}\label{eq:thqg-entanglement-entropy-higher}
S_{\mathrm{EE}}^{(g)}(\cA)
\;:=\;
-\Tr\bigl(\rho_g(\cA) \log \rho_g(\cA)\bigr).
\end{equation}
This satisfies:
\begin{enumerate}[label=\textup{(\alph*)}]
\item $0 \le S_{\mathrm{EE}}^{(g)}(\cA) \le \log E_g(\cA)
  \le \log \dim Q_g(\cA)$.
\item $S_{\mathrm{EE}}^{(1)}(\cA) = 0$
  \textup{(}recovering
  Proposition~\textup{\ref{prop:thqg-III-entanglement-entropy})}.
\item $S_{\mathrm{EE}}^{(g)}(\cA) = S_{\mathrm{EE}}^{(g)}(\cA^!)$
  \textup{(}from complementarity\textup{)}.
\item The shadow tower gives a convergent approximation:
  $S_{\mathrm{EE}}^{(g, \le r)}(\cA) \to
  S_{\mathrm{EE}}^{(g)}(\cA)$ as $r \to \infty$.
\end{enumerate}
\end{proposition}

\begin{proof}[Argument \textup{(}\ClaimStatusHeuristic{}\textup{)}]
Part (a) is the standard von Neumann entropy bound.
Part (b): at genus~$1$, $\rho_1$ has rank~$1$
(one-dimensional $Q_1$), so $S_{\mathrm{EE}}^{(1)} = 0$.
Part (c): the Verdier involution exchanges the roles of
$Q_g(\cA)$ and $Q_g(\cA^!)$, hence exchanges the two
reduced density matrices without changing the spectrum.
Part (d): the shadow tower truncation $\Theta^{\le r}$ converges
to $\Theta$ (Theorem~\ref{thm:recursive-existence}),
and von Neumann entropy is continuous in the trace norm.

The heuristic status reflects the passage from the algebraic
Lagrangian decomposition to the Hilbert space tensor product:
this requires the unitarity inner product, which is additional
data beyond the algebraic framework.  For unitary CFTs (which
include all standard families at generic level), this data exists
and is unique.
\end{proof}

% ======================================================================
%
%  3. THE ALGEBRAIC QUANTUM EXTREMAL SURFACE
%
% ======================================================================

\subsection{The algebraic quantum extremal surface}
\label{subsec:thqg-entanglement-qes}
\index{quantum extremal surface!algebraic|textbf}
\index{QES|see{quantum extremal surface}}

The quantum extremal surface (QES) formula
\cite{Engelhardt-Wall15}
\begin{equation}\label{eq:thqg-qes-physical}
S(R) \;=\; \min_{\gamma}\,\operatorname{ext}\Bigl[
\frac{\mathrm{Area}(\gamma)}{4G_N}
\;+\; S_{\mathrm{bulk}}(\Sigma_\gamma)
\Bigr]
\end{equation}
generalizes the Ryu--Takayanagi formula \cite{Ryu-Takayanagi06}
by including bulk quantum corrections.  The shadow tower provides
the algebraic realization.

\begin{definition}[QES approximant]
\label{def:thqg-qes-approximant}
\index{QES approximant|textbf}
For a gravitational input and integer $r \ge 2$, the
\emph{$r$-th order QES approximant} is:
\begin{equation}\label{eq:thqg-qes-approximant}
S_{\mathrm{QES}}^{(g, \le r)}(\cA)
\;:=\;
\underbrace{
\kappa(\cA) \cdot \lambda_g^{\mathrm{FP}}
}_{\text{area term (arity $2$)}}
\;+\;
\underbrace{
\sum_{j=3}^{r} \delta S^{(g)}_j(\cA)
}_{\text{bulk corrections (arities $3, \ldots, r$)}}
\end{equation}
where $\delta S^{(g)}_j(\cA)$ is the genus-$g$ entanglement
increment contributed by the arity-$j$ shadow jet
$\Sh_j(\cA)$.  At $r = 2$, this is the
classical Ryu--Takayanagi term.
\end{definition}

\begin{theorem}[QES convergence; \ClaimStatusProvedHere]
\label{thm:thqg-qes-convergence}
\index{QES!convergence}
The QES approximants converge:
\[
S_{\mathrm{QES}}^{(g)}(\cA)
\;:=\;
\lim_{r \to \infty}
S_{\mathrm{QES}}^{(g, \le r)}(\cA)
\]
exists for every genus $g \ge 1$ and every gravitational input.
For shadow depth $r_{\max} < \infty$, the series terminates.
For $r_{\max} = \infty$, the convergence rate is controlled
by the shadow growth rate:
$|\delta S^{(g)}_j(\cA)| = O(\rho(\cA)^j \cdot j^{-5/2})$.
\end{theorem}

\begin{proof}
Convergence of the shadow tower
(Theorem~\ref{thm:recursive-existence}) gives convergence of
$\Theta_\cA^{\le r} \to \Theta_\cA$.  The QES approximant
is a continuous functional of $\Theta_\cA^{\le r}$
(it is the genus-$g$ component of the MC element projected
onto the Lagrangian summands), so convergence of the approximants
follows.  The decay estimate is the shadow growth theorem
(Definition~\ref{def:shadow-growth-rate},
Theorem~\ref{thm:shadow-radius}).
\end{proof}

\begin{remark}[The MC equation as extremization]
\label{rem:thqg-mc-as-extremization}
\index{Maurer--Cartan equation!as extremization}
The QES formula \eqref{eq:thqg-qes-physical} requires
\emph{extremization}: the entropy is computed at the extremal surface.
In the algebraic framework, the extremization is the
\emph{Maurer--Cartan equation}:
\[
D_\cA\,\Theta_\cA
+ \tfrac{1}{2}[\Theta_\cA, \Theta_\cA]
+ \hbar\,\Delta(\Theta_\cA) = 0.
\]
The MC element $\Theta_\cA = D_\cA - d_0$ \emph{automatically}
satisfies this equation (Theorem~\ref{thm:mc2-bar-intrinsic}):
the bar-intrinsic construction produces the extremal point
directly, without requiring a variational principle over all
configurations.  This is the algebraic shadow of the QES
being determined by the equations of motion.
\end{remark}

\begin{conjecture}[Spatial entanglement from the open/closed sector]
\label{conj:thqg-spatial-entanglement}
\index{entanglement!spatial|textbf}
\index{Ryu--Takayanagi formula!algebraic|textbf}
Let $\cM$ be an $\cA$-module category at a puncture $p$ of the
tangential log curve $(X, D, \tau)$.  The
\emph{spatial entanglement entropy of the subregion $\cM$} is:
\begin{equation}\label{eq:thqg-spatial-entanglement}
S(\cM) \;=\;
-\Tr_{HH_*(\cM)}\bigl(
\rho_{\cM} \log \rho_{\cM}
\bigr)
\end{equation}
where $\rho_{\cM}$ is the reduced density matrix obtained from
the open/closed MC element $\Theta^{\mathrm{oc}}_\cA$
\textup{(}Construction~\textup{\ref{constr:thqg-oc-mc-element})}
by tracing over the complementary modules.

The annulus trace theorem \textup{(}Theorem~\textup{\ref{thm:thqg-annulus-trace})}
gives the leading contribution:
\[
S(\cM)
\;\sim\;
\log \dim HH_0(\cM)
\;+\;
\text{\textup{(}corrections from $\Theta_\cA^{\le r}$\textup{)}}.
\]
The corrections are controlled by the shadow tower: each arity-$r$
jet of $\Theta_\cA$ contributes a finite correction to $S(\cM)$,
and the total converges by the shadow tower convergence theorem.
\end{conjecture}

% ======================================================================
%
%  4. REPLICA SYMMETRY AND R\'ENYI ENTROPY
%
% ======================================================================

\subsection{Replica symmetry and R\'enyi entropy}
\label{subsec:thqg-entanglement-replica}
\index{replica symmetry|textbf}
\index{R\'enyi entropy|textbf}

\begin{construction}[Replica manifold]
\label{constr:thqg-replica-manifold}
\index{replica manifold|textbf}
For a genus-$g$ surface $\Sigma_g$ and integer $n \ge 2$,
the \emph{$n$-fold replica manifold} is the cyclic $n$-fold
branched cover
$\Sigma_{g,n} \to \Sigma_g$
branched over two points $\{p, q\} \subset \Sigma_g$.
By the Riemann--Hurwitz formula:
\begin{equation}\label{eq:thqg-replica-genus}
g(\Sigma_{g,n}) \;=\; n(g - 1) + 1.
\end{equation}
For $g = 1$: $g(\Sigma_{1,n}) = 1$ (the $n$-fold cover of a
torus branched at two points is again a torus).
For $g = 0$: $g(\Sigma_{0,n}) = 1 - n$ is negative, so the
construction requires $g \ge 1$ for the branched cover to be
a genuine Riemann surface.
\end{construction}

\begin{proposition}[Replica partition function; \ClaimStatusHeuristic]
\label{prop:thqg-replica}
\index{replica!partition function}
The $n$-fold replica partition function of a gravitational input
$(\cA, X, \langle \cdot, \cdot \rangle, k)$ at genus $g \ge 1$
is:
\begin{equation}\label{eq:thqg-replica-pf}
Z_n^{(g)}(\cA)
\;:=\;
\langle \Theta_\cA^{\otimes ?\,} \rangle_{\Sigma_{g,n}}
\end{equation}
\textup{(}the evaluation of the MC element on the replica manifold\textup{)}.
The R\'enyi entropy is:
\begin{equation}\label{eq:thqg-renyi}
S_n^{(g)}(\cA)
\;:=\;
\frac{1}{1-n}
\log\frac{Z_n^{(g)}(\cA)}{Z_1^{(g)}(\cA)^n}.
\end{equation}
The von~Neumann entropy is recovered in the $n \to 1$ limit:
$S^{(g)}(\cA) = \lim_{n \to 1} S_n^{(g)}(\cA)$.
\end{proposition}

\begin{proof}[Argument \textup{(}\ClaimStatusHeuristic{}\textup{)}]
Standard replica argument \cite{Calabrese-Cardy04, Lewkowycz-Maldacena13}
adapted to the algebraic framework.  The
heuristic status reflects two issues: (1) the passage from
integer $n$ to the analytic continuation $n \to 1$ requires
the Carlson theorem, which demands growth bounds on
$Z_n^{(g)}$ that have not been verified in the algebraic setting;
(2) the evaluation on the branched cover $\Sigma_{g,n}$ requires
the sewing construction (MC5) on a surface whose topology depends
on~$n$.
\end{proof}

\begin{remark}[The genus expansion as a replica computation]
\label{rem:thqg-genus-as-replica}
\index{genus expansion!as replica}
The genus expansion
$Z(\cA) = \sum_{g \ge 0} \hbar^{2g-2} Z_g(\cA)$
and the replica computation
$\Tr(\rho^n) = Z_n / Z_1^n$
are related by the Lewkowycz--Maldacena observation:
the saddle point of the gravitational replica path integral
at integer $n$ has genus $g(n)$, which is an increasing function of~$n$.
In the algebraic framework, the genus expansion of the shadow tower
on the replica manifold decomposes as:
\[
Z_n^{(g)}(\cA)
\;=\;
\sum_{h \ge 0} \hbar^{2h-2}\, Z_h(\cA; \Sigma_{g,n}),
\]
where $Z_h(\cA; \Sigma_{g,n})$ is the genus-$h$ partition
function on the replica surface of genus $g(\Sigma_{g,n}) = n(g-1)+1$.
The shadow tower controls each term: $Z_h(\cA; \Sigma_{g,n})$
receives contributions from arities $2, 3, \ldots, r_{\max}$
of the shadow tower, with the scalar term
$\kappa(\cA) \cdot \lambda_h^{\mathrm{FP}}$ giving the
classical R\'enyi entropy.
\end{remark}

% ======================================================================
%
%  5. MODULAR FLOW FROM THE SHADOW CONNECTION
%
% ======================================================================

\subsection{Modular flow from the shadow connection}
\label{subsec:thqg-entanglement-modular-flow}
\index{modular flow|textbf}
\index{Tomita--Takesaki theory!algebraic incarnation|textbf}

\begin{proposition}[Tomita--Takesaki from Verdier;
\ClaimStatusProvedHere]
\label{prop:thqg-tomita-from-verdier}
\index{Verdier involution!as Tomita $J$}
The Verdier involution $\sigma$ on $\cH_g(\cA)$
satisfies:
\begin{enumerate}[label=\textup{(\roman*)}]
\item $\sigma^2 = \id$;
\item $\sigma$ exchanges the two Lagrangian summands:
  $\sigma(Q_g(\cA)) = Q_g(\cA^!)$;
\item $\sigma$ is anti-linear with respect to any
  unitarity inner product compatible with the Verdier pairing;
\item the pairing $\langle v, \sigma(w) \rangle_g$
  is positive-definite on $Q_g(\cA)$
  \textup{(}for unitary $\cA$\textup{)}.
\end{enumerate}
These are the four defining properties of the Tomita conjugation
operator $J$ in the Tomita--Takesaki theory of von~Neumann
algebras.
\end{proposition}

\begin{proof}
Parts (i)--(ii) are the proved properties of the Verdier
involution
(Theorem~\ref{thm:thqg-g3-polarization},
Proposition~\ref{prop:thqg-III-eigenspace}).
Part~(iii): for a unitary chiral algebra, the Shapovalov
form provides a Hermitian structure on the state space.
The Verdier involution, which exchanges $\cA$ and $\cA^!$,
acts as an anti-unitary operator in this Hermitian
structure because Koszul duality reverses the conformal
grading.  Part (iv) follows from the nondegeneracy of
the Verdier pairing restricted to $Q_g \times Q_g^!$
combined with the positivity of the Shapovalov form.
\end{proof}

\begin{definition}[Shadow modular Hamiltonian]
\label{def:thqg-shadow-modular-hamiltonian}
\index{modular Hamiltonian!from shadow metric|textbf}
The \emph{shadow modular Hamiltonian} is the operator
$K_\cA$ on $Q_g(\cA)$ defined by:
\begin{equation}\label{eq:thqg-modular-hamiltonian}
K_\cA \;:=\; -\log \rho_g(\cA)
\end{equation}
where $\rho_g(\cA)$ is the reduced density matrix
\eqref{eq:thqg-density-matrix}.  In the algebraic framework,
$K_\cA$ is determined by the shadow metric
$Q_L$ (Definition~\ref{def:shadow-metric}):
the eigenvalues of $K_\cA$ are the logarithms of the
eigenvalues of $\rho_g$, which are in turn controlled by
the shadow tower truncations $\Theta_\cA^{\le r}$.
\end{definition}

\begin{conjecture}[Shadow connection generates modular flow]
\label{conj:thqg-modular-flow}
\index{modular flow!from shadow connection}
\index{shadow connection!as modular flow generator}
The shadow connection
$\nabla^{\mathrm{sh}} = d - Q_L'/(2Q_L)\, dt$
\textup{(}Theorem~\textup{\ref{thm:shadow-connection})}
generates the modular automorphism group:
the parallel transport of $\nabla^{\mathrm{sh}}$
along the deformation parameter $t$ implements the
modular flow $\sigma_t(a) = \Delta^{it}\, a\, \Delta^{-it}$
of the Tomita--Takesaki theory, where $\Delta = e^{-K_\cA}$
is the modular operator.  In particular:
\begin{enumerate}[label=\textup{(\roman*)}]
\item The logarithmic singularities of $\nabla^{\mathrm{sh}}$
  at the zeros of $Q_L$ correspond to the singular points of
  the modular flow \textup{(}the entanglement temperature
  diverges at these points\textup{)}.
\item The monodromy $-1$ of $\nabla^{\mathrm{sh}}$
  \textup{(}the Koszul sign\textup{)} corresponds to the
  anti-periodicity of the modular flow at inverse temperature
  $\beta = 2\pi$.
\item The shadow depth $r_{\max}(\cA)$ controls the complexity
  of the modular Hamiltonian: for class~G, $K_\cA$ is quadratic;
  for class~L, cubic; and so on.
\end{enumerate}
\end{conjecture}

\begin{remark}[KMS condition from complementarity]
\label{rem:thqg-kms-from-complementarity}
\index{KMS condition!from complementarity}
The KMS condition
$\omega(a\, \sigma_{i\beta}(b)) = \omega(b\, a)$
for a state $\omega$ at inverse temperature $\beta$
is the Tomita--Takesaki analogue of the complementarity constraint
$\Phi(\cA) + \Phi(\cA^!) = \Phi_{\mathrm{total}}$
(Corollary~\ref{cor:thqg-III-complementarity-exchange}).
Both express a symmetry between two halves of a decomposed system:
the KMS condition relates the two orderings of operators in the
thermal state; complementarity relates the two Lagrangian summands
of the gravitational phase space.  The shadow connection provides
the bridge: the parallel transport from $Q_g(\cA)$ to $Q_g(\cA^!)$
is the modular conjugation $J$, and the periodicity $\beta = 2\pi$
is the monodromy of $\nabla^{\mathrm{sh}}$.
\end{remark}

% ======================================================================
%
%  6. ERROR CORRECTION FROM KOSZUL DUALITY
%
% ======================================================================

\subsection{Error correction from Koszul duality}
\label{subsec:thqg-entanglement-error-correction}
\index{error correction!from Koszul duality|textbf}
\index{holographic error correction|textbf}

The bar-cobar adjunction provides a natural encoding/decoding
structure that parallels the quantum error-correcting codes of
holographic entanglement
\cite{Almheiri-Dong-Harlow15, Pastawski-Yoshida-Harlow-Preskill15}.

\begin{proposition}[Bar-cobar as error-correcting code;
\ClaimStatusProvedHere]
\label{prop:thqg-barcobar-error-correction}
\index{bar-cobar!as error-correcting code}
The bar-cobar adjunction $B \dashv \Omega$ endows each
Koszul chiral algebra $\cA$ with the following code structure:
\begin{enumerate}[label=\textup{(\roman*)}]
\item \emph{Encoding.}
  The bar construction $B \colon \cA \to B(\cA)$ maps the
  ``logical'' algebra $\cA$ into the ``physical'' bar
  coalgebra $B(\cA)$.

\item \emph{Decoding.}
  The cobar construction $\Omega \colon B(\cA) \to \Omega(B(\cA))$
  provides the recovery map.

\item \emph{Code perfection on the Koszul locus.}
  The bar-cobar inversion theorem \textup{(}Theorem~B\textup{)}
  gives $\Omega(B(\cA)) \xrightarrow{\sim} \cA$:
  the recovery map is a quasi-isomorphism, recovering $\cA$
  perfectly from the encoded data $B(\cA)$.

\item \emph{Dual code.}
  The Koszul dual $\cA^! = (H^*(B(\cA)))^\vee$ is the
  ``code algebra'': the cohomology of the bar complex, dualized.
  The line-operator category
  $\mathcal{C} \simeq \cA^!\text{-}\mathsf{mod}$
  is the category of ``codewords.''
\end{enumerate}
\end{proposition}

\begin{proof}
The encoding~(i) and decoding~(ii) are the bar and cobar functors,
which are defined for any augmented dg algebra/coalgebra.
Code perfection~(iii) is Theorem~B (bar-cobar inversion on the
Koszul locus).  The dual code~(iv) is the definition of the
Koszul dual and the identification of the line-operator category
(Theorem~\ref{thm:thqg-g5-yangian}).
\end{proof}

\begin{conjecture}[Shadow depth as code distance]
\label{conj:thqg-shadow-depth-code-distance}
\index{shadow depth!as code distance}
\index{code distance|textbf}
The shadow depth $r_{\max}(\cA)$ controls the error-correcting
capacity of the holographic code:
\begin{center}
\renewcommand{\arraystretch}{1.3}
\small
\begin{tabular}{clll}
\textbf{Class} & $r_{\max}$ & \textbf{Code property}
  & \textbf{Physical meaning} \\
\midrule
G & $2$ & trivial code \textup{(}classical\textup{)}
  & Gaussian/free theory \\
L & $3$ & single-error correction
  & tree-level interactions \\
C & $4$ & double-error correction
  & contact/quartic interactions \\
M & $\infty$ & unbounded capacity
  & fully interacting QFT
\end{tabular}
\end{center}
The interpretation: the shadow tower
$\Theta_\cA^{\le r}$ encodes gravitational data
at arity $r$.  Erasure of the arity-$r$ data can be
recovered from arities $\le r-1$ if and only if the
$(r+1)$-th obstruction class $o^{(r+1)}(\cA)$ vanishes ---
precisely when $r \le r_{\max}$.
\end{conjecture}

\begin{theorem}[Entanglement wedge from the open/closed realization;
\ClaimStatusProvedHere]
\label{thm:thqg-entanglement-wedge}
\index{entanglement wedge|textbf}
\index{subregion duality!algebraic|textbf}
The open/closed realization
\textup{(}\S\textup{\ref{sec:thqg-open-closed-realization})}
provides a subregion decomposition of the holographic system:
\begin{enumerate}[label=\textup{(\roman*)}]
\item \emph{Subregion.}
  A boundary module $\cM$ at a puncture~$p$ defines the
  ``subregion'' --- the open sector of the Swiss-cheese
  decomposition
  (Theorem~\textup{\ref{thm:thqg-swiss-cheese}}).

\item \emph{Entanglement wedge.}
  The bulk operators reconstructible from the subregion $\cM$
  are precisely those in the image of the open/closed map
  $\cO \colon C^\bullet_{\mathrm{ch}}(\cA, \cA) \to
  \operatorname{End}(\cM)$:
  a bulk operator $\Phi \in C^\bullet_{\mathrm{ch}}(\cA, \cA)$
  is in the entanglement wedge of $\cM$ if and only if
  $\cO(\Phi) \ne 0$.

\item \emph{Bulk reconstruction.}
  By the open/closed projection principle
  \textup{(}Theorem~\textup{\ref{thm:thqg-oc-projection})},
  the full holographic datum $\cH(\cA)$ is recoverable from
  the open/closed MC element $\Theta^{\mathrm{oc}}_\cA$.
  The closed projection $\pi_{\mathrm{cl}}$ recovers
  $\Theta_\cA$ (the full bulk data), while the open projection
  recovers the boundary module structures.

\item \emph{Partial trace.}
  The annulus trace theorem
  \textup{(}Theorem~\textup{\ref{thm:thqg-annulus-trace})}
  computes $HH_*(\cM)$ as the factorization homology over
  the boundary circle.  This is the algebraic partial trace:
  it integrates out the boundary degrees of freedom at
  the puncture, leaving the Hochschild invariants.
\end{enumerate}
\end{theorem}

\begin{proof}
All four parts follow from the proved open/closed realization:
(i)~is the Swiss-cheese theorem;
(ii)~is the definition of the open/closed map combined with the
chiral derived center
(Definition~\ref{def:thqg-chiral-derived-center});
(iii)~is the open/closed projection principle;
(iv)~is the annulus trace theorem.
\end{proof}

% ======================================================================
%
%  7. THE ALGEBRAIC PAGE CONSTRAINT
%
% ======================================================================

\subsection{The algebraic Page constraint}
\label{subsec:thqg-entanglement-page}
\index{Page constraint!algebraic|textbf}
\index{Page curve|textbf}
\index{complementarity!as Page constraint}

\begin{theorem}[Algebraic Page constraint; \ClaimStatusProvedHere]
\label{thm:thqg-page-constraint}
For a gravitational input $(\cA, X, \langle \cdot, \cdot \rangle, k)$
and any genus $g \ge 1$:
\begin{enumerate}[label=\textup{(\roman*)}]
\item \textup{(Page bound.)}
  \begin{equation}\label{eq:thqg-page-bound}
  S_{\mathrm{EE}}^{(g)}(\cA)
  \;\le\;
  \min\bigl(\log \dim Q_g(\cA),\;
  \log \dim Q_g(\cA^!)\bigr).
  \end{equation}

\item \textup{(Complementarity sum.)}
  \begin{equation}\label{eq:thqg-page-complementarity}
  \dim Q_g(\cA) + \dim Q_g(\cA^!)
  \;=\;
  \dim \cH_g(\cA)
  \;=\;
  \dim H^*(\overline{\cM}_g,\, \cZ(\cA)).
  \end{equation}

\item \textup{(Self-dual saturation.)}
  For $\cA \cong \cA^!$, the Page bound is symmetric:
  $S_{\mathrm{EE}}^{(g)} \le \frac{1}{2}
  \log \dim \cH_g(\cA)$.

\item \textup{(Scalar complementarity.)}
  At the scalar level, the Page constraint for the
  Virasoro family reads:
  \begin{equation}\label{eq:thqg-page-virasoro-scalar}
  F_g^{\mathrm{sc}}(\mathrm{Vir}_c)
  + F_g^{\mathrm{sc}}(\mathrm{Vir}_{26-c})
  \;=\;
  13 \cdot \lambda_g^{\mathrm{FP}},
  \end{equation}
  which is the scalar complementarity identity of Theorem~D.
\end{enumerate}
\end{theorem}

\begin{proof}
Part (i): the entropy bound is the standard von Neumann bound
applied to the reduced density matrix on the smaller factor.
Part (ii): the complementarity theorem
(Theorem~\ref{thm:quantum-complementarity-main}).
Part (iii): Corollary~\ref{cor:thqg-III-dimension-parity}.
Part (iv): linearity of $F_g^{\mathrm{sc}}$ in $\kappa$
plus $\kappa(\mathrm{Vir}_c) + \kappa(\mathrm{Vir}_{26-c})
= c/2 + (26-c)/2 = 13$.
\end{proof}

\begin{remark}[The Page transition at the self-dual point]
\label{rem:thqg-page-transition}
\index{Page transition|textbf}
\index{self-dual point!as Page transition}
For the Virasoro family parametrized by central charge $c$:
\begin{itemize}
\item $c > 13$: $\kappa(\mathrm{Vir}_c) > \kappa(\mathrm{Vir}_{26-c})$,
  so the boundary sector dominates at the scalar level.
  The gravitational entropy is controlled by $Q_g(\cA^!) \sim
  (26-c)/2$: this is the \emph{Hawking phase}, where the bulk
  contribution is smaller.
\item $c = 13$: $\kappa = \kappa'$, the self-dual point.
  The boundary and bulk contributions are equal.  This is the
  algebraic incarnation of the \emph{Page time}: the transition
  between the Hawking phase and the information-recovery phase.
\item $c < 13$: $\kappa(\mathrm{Vir}_c) < \kappa(\mathrm{Vir}_{26-c})$,
  and the bulk sector dominates.  The entanglement entropy is
  controlled by $Q_g(\cA) \sim c/2$: this is the
  \emph{information-recovery phase}.
\end{itemize}
The Virasoro Koszul involution $c \mapsto 26 - c$ acts as the
Page reflection: it interchanges the two phases.  The fixed point
$c = 13$ is the unique self-dual Virasoro algebra
(Proposition~\ref{prop:thqg-V-c13-self-duality}).
\end{remark}

\begin{remark}[Page curve from shadow tower convergence]
\label{rem:thqg-page-curve}
\index{Page curve!from shadow tower}
The shadow tower provides a \emph{constructive, convergent}
approximation to the Page constraint at each order:
\begin{enumerate}[label=\textup{(\arabic*)}]
\item At arity $2$ (scalar level): the Page constraint is the
  scalar complementarity \eqref{eq:thqg-page-virasoro-scalar}.
\item At arity $3$ (cubic shadow): the cubic obstruction $C(\cA)$
  and $C(\cA^!)$ satisfy a complementarity relation
  (independent sum factorization,
  Proposition~\ref{prop:independent-sum-factorization}).
\item At arity $4$ (quartic shadow): the quartic resonance
  classes $R_4^{\mathrm{mod}}(\cA)$ and $R_4^{\mathrm{mod}}(\cA^!)$
  are constrained by the discriminant complementarity
  $\Delta(\cA) + \Delta(\cA^!) = \text{constant}$
  (Theorem~\ref{thm:shadow-connection}).
\item At all arities: the full shadow tower
  $\Theta_\cA = \varprojlim \Theta_\cA^{\le r}$
  satisfies the complete Page constraint at all genera simultaneously.
\end{enumerate}
This is the unique property of the programme: no other approach to
quantum gravity provides a convergent, finite-order approximation
scheme for the entanglement constraints.
\end{remark}

\begin{proposition}[Entanglement capacity from shadow depth;
\ClaimStatusProvedHere]
\label{prop:thqg-entanglement-capacity}
\index{entanglement capacity|textbf}
The \emph{entanglement capacity} at genus $g$ is
\begin{equation}\label{eq:thqg-entanglement-capacity}
C_g(\cA)
\;:=\;
\log \dim Q_g(\cA).
\end{equation}
For the Virasoro family at genus $g \ge 2$:
\begin{enumerate}[label=\textup{(\alph*)}]
\item At the scalar level:
  $C_g^{\mathrm{sc}}(\mathrm{Vir}_c) = \log 1 = 0$
  \textup{(}the scalar contribution is one-dimensional\textup{)}.
\item Including the genus-$1$ Hessian correction
  $\delta H^{(1)}_{\mathrm{Vir}} = 120/[c^2(5c+22)] \cdot x^2$:
  the entanglement capacity receives its first nontrivial
  contribution from the quartic shadow.
\item Complementarity constrains the total capacity:
  $C_g(\cA) + C_g(\cA^!) \le
  \log \dim \cH_g(\cA)$.
\end{enumerate}
\end{proposition}

\begin{proof}
Part~(a): the scalar MC element
$\Theta^{(g),\le 2} = \kappa \cdot \lambda_g^{\mathrm{FP}}$
is one-dimensional in $Q_g(\cA)$ at every genus.
Part~(b): the quartic shadow
$\mathfrak{Q}^{(0)}_\cA$ contributes to the genus-$1$ Hessian
via the genus loop operator
(Corollary~\ref{cor:genus-1-hessian-from-quartic}),
providing the first tensorial modular invariant that can increase
$\dim Q_g$ above~$1$.
Part~(c): the complementarity sum
(Theorem~\ref{thm:quantum-complementarity-main}).
\end{proof}

% ======================================================================
%
%  8. SIX THEOREM TARGETS
%
% ======================================================================

\subsection{Six theorem targets extending the THQG programme}
\label{subsec:thqg-entanglement-targets}
\index{THQG programme!entanglement extension|textbf}

We formulate six theorem targets (G11--G16) extending the
ten-theorem programme of \S\ref{sec:thqg-ten-theorems}
to the entanglement domain.  Each deploys the shadow tower
on a problem where no other approach has a convergent
approximation scheme.

\begin{target}[\textbf{G11}: Algebraic entanglement entropy]
\label{target:thqg-g11}
\index{G11|textbf}
\index{entanglement entropy!algebraic}
For a gravitational input, the genus-$g$ entanglement entropy
$S_{\mathrm{EE}}^{(g)}(\cA)$ is well-defined, satisfies the
Page bound \eqref{eq:thqg-page-bound}, and admits a convergent
shadow tower expansion:
\begin{equation}\label{eq:thqg-g11}
S_{\mathrm{EE}}^{(g)}(\cA)
\;=\;
\sum_{r=2}^{r_{\max}} \delta S_r^{(g)}(\cA)
\end{equation}
with $\delta S_2^{(g)} = 0$ \textup{(}scalar term contributes
no entanglement\textup{)} and each $\delta S_r^{(g)}$
computable from the shadow jet $\Sh_r(\cA)$.
\medskip

\noindent
\textbf{Status:} Proved at genus $1$
(Proposition~\ref{prop:thqg-III-entanglement-entropy}).
Heuristic at genus $g \ge 2$
(Proposition~\ref{prop:thqg-entanglement-entropy-higher}).
The proof gap is the passage from the algebraic Lagrangian
decomposition to the Hilbert space tensor product
(requires unitarity).
\end{target}

\begin{target}[\textbf{G12}: Algebraic quantum extremal surface]
\label{target:thqg-g12}
\index{G12|textbf}
\index{quantum extremal surface!algebraic}
The spatial entanglement entropy of a boundary subregion $\cM$
is determined by the open/closed MC element
$\Theta^{\mathrm{oc}}_\cA$
(Conjecture~\ref{conj:thqg-spatial-entanglement}):
\[
S(\cM) \;=\; \kappa(\cA) \cdot \lambda_g^{\mathrm{FP}}
\;+\; \sum_{r \ge 3}
\delta S_r^{\mathrm{oc}}(\cM; \Theta_\cA).
\]
The first term is the algebraic Ryu--Takayanagi area; the
corrections are the QES bulk entropy.
The MC equation is the extremization condition.
\medskip

\noindent
\textbf{Status:} Conjectured
(Conjecture~\ref{conj:thqg-spatial-entanglement}).
Convergence of the approximants is proved
(Theorem~\ref{thm:thqg-qes-convergence}).
The identification with the physical QES formula requires
matching the open/closed MC decomposition with the
Engelhardt--Wall prescription.
\end{target}

\begin{target}[\textbf{G13}: Modular flow from the shadow connection]
\label{target:thqg-g13}
\index{G13|textbf}
\index{modular flow!from shadow connection}
The shadow connection $\nabla^{\mathrm{sh}}$ generates
the modular automorphism group of the Tomita--Takesaki
theory (Conjecture~\ref{conj:thqg-modular-flow}):
the Verdier involution $\sigma$ is the modular conjugation $J$,
the shadow metric $Q_L$ determines the modular Hamiltonian
$K_\cA$, and the KMS condition at $\beta = 2\pi$ follows
from the monodromy of $\nabla^{\mathrm{sh}}$.
\medskip

\noindent
\textbf{Status:} The structural identification
$\sigma = J$ is proved
(Proposition~\ref{prop:thqg-tomita-from-verdier}).
The full modular flow identification is conjectured.
\end{target}

\begin{target}[\textbf{G14}: Holographic error correction]
\label{target:thqg-g14}
\index{G14|textbf}
\index{holographic error correction}
The bar-cobar adjunction provides a quantum error-correcting
code structure for the holographic system
(Proposition~\ref{prop:thqg-barcobar-error-correction}):
encoding via $B$, decoding via $\Omega$, code perfection
from Koszulness, entanglement wedge from the open/closed
realization
(Theorem~\ref{thm:thqg-entanglement-wedge}).
The shadow depth $r_{\max}(\cA)$ classifies the code complexity
(Conjecture~\ref{conj:thqg-shadow-depth-code-distance}).
\medskip

\noindent
\textbf{Status:} The code structure is proved
(Proposition~\ref{prop:thqg-barcobar-error-correction},
Theorem~\ref{thm:thqg-entanglement-wedge}).
The shadow depth as code distance is conjectured.
\end{target}

\begin{target}[\textbf{G15}: Algebraic Page constraint]
\label{target:thqg-g15}
\index{G15|textbf}
\index{Page constraint}
The complementarity theorem imposes a Page-type bound on the
entanglement entropy at all genera
(Theorem~\ref{thm:thqg-page-constraint}):
\[
S_{\mathrm{EE}}^{(g)}(\cA)
\;\le\;
\min\bigl(C_g(\cA), C_g(\cA^!)\bigr),
\qquad
C_g(\cA) + C_g(\cA^!) = \log \dim \cH_g(\cA).
\]
The shadow tower provides a convergent finite-order approximation
to both sides of the bound.  The self-dual point
$\cA \cong \cA^!$ saturates the constraint
(Remark~\ref{rem:thqg-page-transition}).
\medskip

\noindent
\textbf{Status:} Proved
(Theorem~\ref{thm:thqg-page-constraint}).
The identification with the physical Page curve of black hole
evaporation is heuristic.
\end{target}

\begin{target}[\textbf{G16}: Replica structure]
\label{target:thqg-g16}
\index{G16|textbf}
\index{replica structure}
The R\'enyi entropy $S_n^{(g)}(\cA)$ is computed by the
shadow tower on the replica manifold $\Sigma_{g,n}$
(Proposition~\ref{prop:thqg-replica}):
\[
S_n^{(g)}(\cA)
\;=\;
\frac{1}{1-n}
\log\frac{Z_{n(g-1)+1}(\cA)}{Z_g(\cA)^n}.
\]
The genus expansion of each term is controlled by the shadow tower.
The von~Neumann entropy is the $n \to 1$ limit.
\medskip

\noindent
\textbf{Status:} Heuristic
(Proposition~\ref{prop:thqg-replica}).
The analytic continuation $n \to 1$ requires growth bounds
not yet established.
\end{target}

\begin{remark}[Dependency structure of G11--G16]
\label{rem:thqg-g11-g16-dependencies}
\index{THQG programme!dependency structure}
The six targets depend on the original ten as follows:
\begin{center}
\begin{tikzpicture}[
  node distance=1.2cm and 1.8cm,
  every node/.style={draw, rounded corners, inner sep=4pt,
    font=\small},
  proved/.style={fill=green!15},
  conjectural/.style={fill=yellow!15},
  heuristic/.style={fill=orange!15}
]
\node[proved] (g3) {G3};
\node[proved, right=of g3] (g4) {G4};
\node[proved, right=of g4] (g10) {G10};
\node[heuristic, below left=of g3] (g11) {G11};
\node[conjectural, below=of g3] (g12) {G12};
\node[conjectural, below right=of g3] (g13) {G13};
\node[proved, below=of g4] (g14) {G14};
\node[proved, below right=of g4] (g15) {G15};
\node[heuristic, below=of g10] (g16) {G16};
\draw[->] (g3) -- (g11);
\draw[->] (g3) -- (g12);
\draw[->] (g3) -- (g13);
\draw[->] (g3) -- (g15);
\draw[->] (g4) -- (g15);
\draw[->] (g10) -- (g16);
\draw[->] (g11) -- (g12);
\end{tikzpicture}
\end{center}
\textup{(}Green = proved, yellow = conjectural,
orange = heuristic.\textup{)}
G3 (symplectic polarization) and G4 ($S$-duality) are the primary
inputs; G11 (entanglement entropy) feeds into G12 (QES).
\end{remark}

\begin{remark}[The convergent advantage]
\label{rem:thqg-convergent-advantage}
\index{convergence!advantage over other programmes}
The unique structural advantage of the shadow tower programme over
competing approaches:
\begin{center}
\renewcommand{\arraystretch}{1.3}
\small
\begin{tabular}{lll}
\textbf{Programme} & \textbf{Expansion} & \textbf{Convergence} \\
\midrule
String theory genus expansion & $\sum_g g_s^{2g-2} F_g$
  & asymptotic ($\sim (2g)!$) \\
Loop QG spin foams & $\sum_\sigma A_\sigma$
  & unknown \\
AdS/CFT large $N$ & $1/N^2$ expansion
  & asymptotic (typically) \\
\textbf{Shadow tower} & $\Theta_\cA^{\le r}$
  & \textbf{proved convergent} ($O(\rho^r r^{-5/2})$)
\end{tabular}
\end{center}
The shadow tower converges because each arity-$r$ contribution
is controlled by the shadow growth rate $\rho(\cA) < 1$
(for $c > c^* \approx 6.125$ in the Virasoro family).
For $c < c^*$, the tower is Borel summable with known Stokes
data (Theorem~\ref{thm:shadow-radius}).  In either regime, the
entanglement programme has a systematic, finite-order
approximation scheme with proved error bounds.
\end{remark}
